% ============================================
% ЗАКЛЮЧЕНИЕ
% Подводятся итоги преддипломной практики:
% выполнение задач этапов 2–4 плана-графика, формулировка
% критериев достижимости цели и основных результатов ВКР.
% ============================================

В ходе преддипломной практики я выполнил все задачи,
предусмотренные индивидуальным планом-графиком прохождения
практики, и в установленные сроки подготовил полный комплект
материалов выпускной квалификационной работы к итоговой защите.

В рамках этапа~2 плана-графика подготовлен полный текст ВКР
объёмом~78~страниц основной части (с учётом введения и
заключения), что с запасом удовлетворяет требованию <<не менее
40~страниц для бакалаврской ВКР>>. Текст оформлен в системе
вёрстки \LaTeX{} с использованием шаблона, реализующего все
требования к оформлению, изложенные в методических
рекомендациях ИТМО. Структура работы включает все необходимые
элементы~--- титульный лист, аннотацию, индивидуальное задание,
содержание, список сокращений и условных обозначений, список
терминов и определений, введение, четыре раздела основной части
с выводами по каждому разделу, заключение, список использованных
источников и девять приложений. Введение содержит общую
характеристику работы в соответствии с требованиями:
актуальность, степень разработанности, решаемую проблему, цель и
пять задач, практическую значимость, а также добавленные по
согласованию с научным руководителем блоки <<Апробация и
внедрение результатов>> и <<Личный вклад автора>>.

Список использованных источников сформирован из~29~наименований
и удовлетворяет всем формальным требованиям с запасом: из них
21~источник опубликован не ранее~2021~года (норма <<не
менее~20>>); 7~источников представляют собой научные статьи в
журналах, рекомендованных ВАК~РФ для диссертаций, что
составляет~24,1\% от общего числа источников (норма <<не
менее~20\%>>); 13~источников представляют собой статьи в
изданиях, индексируемых международными базами Web of Science
и Scopus, что составляет~44,8\% от общего числа источников
(норма <<не менее~15\%>>). В составе ВАК-источников
представлены журналы <<Научно-технический вестник
информационных технологий, механики и оптики>>, <<Цифровая
обработка сигналов>>, <<Автоматика и телемеханика>>,
<<Компьютерные исследования и моделирование>>, <<Компьютерная
оптика>>, относящиеся к различным категориям перечня ВАК (от
К1 до К3), что обеспечивает разнообразие источниковой базы
работы.

Отдельного внимания при подготовке текста ВКР заслуживает
выработка критериев достижимости цели, поскольку именно они
определяют, в чём конкретно состоит преимущество предложенного
решения и каким образом это преимущество доказывается. Подход к
формулировке критериев сложился по итогам двух очных консультаций
с научным руководителем. На этих консультациях руководитель
обозначил, что во втором разделе ВКР должно быть представлено
авторское решение, чем-либо превосходящее существующие, и что
для подтверждения этого превосходства необходимо ввести
несколько измеримых метрик, по которым можно судить о выполнении
или невыполнении требований, сформулированных по итогам анализа
предметной области. В качестве отправных примеров метрик
руководитель предложил рассмотреть энергопотребление, дальность
(радиус) работы устройства и <<ёмкость>> нейросетевой
архитектуры~--- размер модели и число её параметров.

Опираясь на эти ориентиры, по итогам анализа предметной области
я сформулировал требования к разрабатываемой системе и впервые
представил соответствующие критерии достижимости на первом
прослушивании. Эта первая редакция критериев оказалась
неудачной: в ней были заданы заведомо недостижимые целевые
значения по энергопотреблению и точности модели, а сами метрики
слабо отражали суть и реальные результаты работы. По итогам
обсуждения критерии были существенно переработаны: ко второму
прослушиванию я конкретизировал их и озвучил устно, получил
правки, и к предзащите сформировал окончательный набор из пяти
критериев~C1--C5, привязанных к измеримым величинам и к
требованиям бенчмарка MLPerf~Tiny:

\begin{itemize}
	\item \textit{C1~--- точность распознавания} не ниже~90\%
	      на тестовой выборке набора Google Speech Commands в
	      соответствии с порогом бенчмарка MLPerf~Tiny;
	\item \textit{C2~--- размер квантизованной модели} не
	      более~128~КБ для размещения во внутреннем SRAM
	      платформы, что необходимо для корректной работы
	      аппаратного SIMD-ускорения;
	\item \textit{C3~--- потеря точности от квантизации} не
	      более~1~процентного пункта относительно
	      FP32-варианта модели;
	\item \textit{C4~--- латентность инференса} не более~500~мс,
	      что соответствует требованиям интерактивного
	      голосового интерфейса;
	\item \textit{C5~--- снижение средней потребляемой мощности}
	      не менее чем в~4~раза относительно режима непрерывного
	      инференса на той же платформе.
\end{itemize}

Именно этот набор критериев представлен в тексте ВКР и на слайде
анализа предметной области презентации; задача экспериментальной
части работы состояла в подтверждении выполнения каждого из них.

Заключение текста ВКР содержит список основных результатов
исследования с количественными и качественными характеристиками,
сгруппированный по пяти задачам и подтверждающий достижение
поставленной цели. Кратко эти результаты состоят в следующем.

По задаче~1 проведён системный анализ предметной области с
рассмотрением облачных решений, специализированных аппаратных
платформ (ASIC, ПЛИС, DSP) и микроконтроллеров общего
назначения; по его итогам сформулированы требования к системе и
критерии достижимости~C1--C5. По задаче~2 разработана каскадная
архитектура детекции из трёх ступеней возрастающей вычислительной
сложности и выведена аналитическая модель средней потребляемой
мощности, из которой получены два проектных критерия~---
кратное разделение ступеней по мгновенному потреблению и по
частоте активации. По задаче~3 проведено систематическое
исследование пространства архитектур DS-CNN (обучено
134~архитектуры, оценено 239~точек квантизации), показавшее, что
98\% парных различий точности методов PTQ и QAT укладываются в
коридор статистической погрешности~$\pm 0{,}65$~п.п. По задаче~4
разработан двухдоменный лабораторный измерительный стенд с
регистрацией мгновенного тока с разрешением единиц микроампер и
выполнена экспериментальная апробация системы на платформе
ESP32-S3. По задаче~5 выполнена покомпонентная декомпозиция
энергобюджета цикла каскадной обработки, выявившая доминирующий
вклад сервисных этапов (инициализация и захват аудио).

Выполнение критериев достижимости цели подтверждено
экспериментально по каждому из пяти критериев со значительным
запасом:

\begin{itemize}
	\item по критерию~C1 достигнута точность распознавания
	      96,12\% при требовании не ниже~90\%;
	\item по критерию~C2 размер выбранной квантизованной
	      модели составил~101~КБ при требовании не более~128~КБ;
	\item по критерию~C3 потеря точности от квантизации
	      отсутствует~--- квантизованная INT8-модель показала
	      точность не ниже FP32-варианта (различие в пределах
	      статистической погрешности), при требовании потери не
	      более~1~процентного пункта;
	\item по критерию~C4 латентность инференса составила~460~мс
	      при требовании не более~500~мс;
	\item по критерию~C5 расчётное снижение средней потребляемой
	      мощности превысило~40~раз при требовании не менее
	      чем в~4~раза.
\end{itemize}

Научная новизна работы, отражённая в заключении ВКР, состоит из
трёх элементов: систематического исследования пространства
архитектур DS-CNN со статистической оценкой значимости различий
по доверительному интервалу Уилсона; покомпонентной декомпозиции
энергобюджета каскадного цикла, впервые выделившей сервисные
этапы как доминирующие компоненты энергобюджета; и сравнительного
анализа методов квантизации PTQ и QAT на большой выборке моделей,
статистически подтвердившего их эквивалентность для квантизации
до INT8. В заключении ВКР приведены также обоснование практической
значимости работы и шесть направлений дальнейших исследований по
теме~--- проектирование собственной печатной платы с LDO
ультранизкого тока покоя, замена аналогового детектора первой
ступени на MEMS-сенсор с пробуждением по звуку, применение
кольцевого буфера на сопроцессоре ULP, сокращение длительности
этапа инициализации, верификация метода на платформах STM32U5 и
nRF52840 и исследование стратегий постобработки выходов
классификатора.

В рамках этапа~3 плана-графика текст ВКР успешно прошёл проверку
в системе <<Антиплагиат.ВУЗ>>~20~мая, проведённую секретарём ГЭК.
Результат проверки~--- оригинальность~97,36\% при минимально
требуемом значении~75\%, то есть с запасом более чем в~22
процентных пункта. Доля совпадений составила~2,11\%,
цитирования~--- 0,53\%, самоцитирование и доля ИИ-контента
нулевые. Высокий уровень оригинальности обусловлен
самостоятельным характером проведённой работы: оригинальной
формулировкой проектных критериев каскадной архитектуры,
самостоятельной реализацией программно-аппаратного комплекса и
лабораторного измерительного стенда, авторской методикой
систематического исследования пространства архитектур и
декомпозиции энергобюджета цикла каскадной обработки. Полученная
справка о результатах проверки приведена в
приложении~\ref{app:antiplagiat}.

В рамках этапа~4 плана-графика подготовлена электронная
презентация результатов работы из~15~слайдов основной части и
финального слайда благодарности, выдержанная в едином графическом
стилевом оформлении (фиолетово-сиреневая палитра, белый фон,
единая шапка с логотипом университета). Презентация подготовлена в
сервисе Google~Slides на основе официального шаблона ИТМО; для
демонстрации работы устройства отдельно подготовлен и встроен в
презентацию короткий видеоролик, в котором показаны включение
лабораторного стенда и две успешные классификации голосовых
команд с синхронной индикацией состояний каскада и графика
мгновенного тока. Презентация для предзащиты в её исходном виде
приведена в приложении~\ref{app:presentation}.

Процедура предзащиты пройдена~21~мая в дистанционном формате.
Получены следующие оценки: готовность ВКР~--- 95\%, готовность
презентации~--- 85\%, общая готовность к защите~--- 80\%. По
итогам выступления получен допуск к итоговой защите. Научный
руководитель направил подробные письменные комментарии по
доработке отдельных слайдов презентации. Доработка презентации
ведётся в составе подготовки к итоговой защите и должна быть
завершена не позднее чем за неделю до её проведения.

Оформление отчётной документации по практике выполнено в
соответствии с требованиями. Структура настоящего отчёта
включает все необходимые элементы: титульный лист, введение,
основную часть с подробным описанием выполнения задач~2--4
плана-графика практики, заключение, список использованных
источников и приложения, в которых размещены справка о
прохождении проверки в системе <<Антиплагиат>>, презентация для
предзащиты, письменные комментарии научного руководителя по
итогам предзащиты и план-график прохождения преддипломной
практики.

По итогам преддипломной практики выпускная квалификационная
работа находится в состоянии полной готовности к итоговой защите.
Все формальные требования к содержанию, оформлению и проверке
оригинальности текста выполнены с запасом по каждому из
показателей, а выполнение всех пяти критериев достижимости цели
подтверждено экспериментально. Доработка презентации по
комментариям научного руководителя и проведение финальных
репетиций выступления планируются завершить в течение недели,
предшествующей дате защиты.
